\begin{donnees}
    \begin{itemize}
        \item $m=\qty{0,24}{\kg}$~; $\pi^{2}=10$~;
              les frottements sont négligeables.
    \end{itemize}
\end{donnees}

On écarte le solide $\mathrm{(S)}$ de la Fig.~\ref{fig:solide_ressort} de sa
position d’équilibre d’une distance $X_{m}$ et on l'abandonne sans vitesse
initiale.

\begin{enumerate}
    \item L'étude expérimentale a permis d'obtenir la courbe de la
          Fig.~\ref{fig:2bac_svt_national_2019_rattrapage_phys_ex3_3}
          représentant les variations de l'accélération $\ddot{x}(t)$ du
          mouvement de $G$.\par
          \begin{minipage}{0.68\linewidth}
              \begin{enumerate}
                  \item En appliquant la deuxième loi de Newton, établir
                        l'équation différentielle vérifiée par l'abscisse $x$.
                  \item La solution de cette équation différentielle s'écrit~:\\
                        $x(t)=X_{m}\cdot
                        \cos\left(\frac{2\pi}{T_{0}}\cdot t\right)$

                         Trouver en fonction des paramètres utiles,
                         l’expression de l’accélération $\ddot{x}(t)$.
                   \item En exploitant la courbe de la
                         Fig.~\ref{fig:2bac_svt_national_2019_rattrapage_phys_ex3_3},
                         déterminer les valeurs de $T_{0}$ et de $X_{m}$.
                   \item Déduire la valeur de la raideur $K$.
              \end{enumerate}
          \end{minipage}
          \hfill
          \begin{minipage}{0.30\linewidth}
              \begin{figure}[H]
                  \centering
                  \begin{tikzpicture}
                      \def\Xm{0.4}
                      \def\T{2}
                      \pgfmathpi\let\pi=\pgfmathresult
                      \begin{axis}[
                              xmin=0,xmax=4.5,
                              ymin=-0.44,ymax=0.59,
                              width=6cm,height=6cm,
                              trig format plots=rad,
                              xtick distance=1, ytick distance=0.2,
                              xticklabels={,,1},
                              yticklabels={,,,,{0,2}},
                              xlabel=$t~(\unit{\second})$,
                              ylabel=$\ddot{x}~(\unit{\a})$,
	                          every axis x label/.append style={
	                            at={(1.03,0.45)},anchor=south east,
	                          },
                              minor tick num=0,
                          ]
                          \addplot[domain=0:4,samples=200,very thick,blue]
                              {\Xm*cos(2*pi/\T*x+pi)};
                      \end{axis}
                  \end{tikzpicture}
                  \caption{}
                  \label{fig:2bac_svt_national_2019_rattrapage_phys_ex3_3}
              \end{figure}
          \end{minipage}
    \item Déterminer dans l'intervalle $[0;\qty{3}{\second}]$ les instants où la
          vitesse de $G$ est maximale. Calculer la valeur de cette vitesse.
\end{enumerate}

